\section{Generalization: Stability, Capacity, and Optimization Error}
\label{app:generalization}

Generalization is a statement about the represented predictor and the data
distribution. It cannot be deduced from a small vertical-energy statistic.
We give two complementary analyses. Algorithmic stability controls the
sensitivity of the actual training procedure to one changed example.
Capacity bounds control a specified class of products, independently of the
training procedure. The former needs update-map regularity, including state;
the latter needs explicit bounds on the represented weights. Balanced factor
norms make these statements invariant to the arbitrary $\mathrm{GL}(r)$
coordinates, but balancing alone does not impose either assumption.

\subsection{Statistical setting and loss sensitivity}
Let $S=(z_1,\ldots,z_n)$ be independent samples from a distribution
$\mathcal D$. For a product parameter $X$ define
\begin{equation}
 R(X)=\mathbb E_{z\sim\mathcal D}\ell(X;z),\qquad
 R_S(X)=\frac1n\sum_{i=1}^n\ell(X;z_i).
\end{equation}
Unless otherwise indicated, assume $0\leq\ell\leq1$ and
\begin{equation}
 |\ell(X;z)-\ell(Y;z)|\leq L_\ell\|X-Y\|.
 \label{eq:loss-product-lipschitz}
\end{equation}
Squared loss or cross-entropy is not globally bounded without restricting
outputs or the parameter domain; the bounded-loss theorems apply to a
bounded or clipped loss, or to a domain on which an explicit bound is
available. Replacing the unit range by $[0,M]$ multiplies the concentration
terms by $M$.

\begin{lemma}[Product sensitivity along the quotient]
\label{lem:product-sensitive}
If a connecting quotient path lies in $\mathcal K_{a,b}$, then
\begin{equation}
 \|X-Y\|\leq\sqrt{2b}\,d_g(X,Y),
 \quad |\ell(X;z)-\ell(Y;z)|
       \leq L_\ell\sqrt{2b}\,d_g(X,Y),
 \label{eq:loss-quotient-lipschitz}
\end{equation}
where $d_g$ is the infimum of lengths of such admissible connecting paths.
\end{lemma}
\begin{proof}
For every path $\gamma$, Eq.~\eqref{eq:metric-equivalence} implies
$\|\dot\gamma\|\leq\sqrt{2b}\|\dot\gamma\|_g$.
Integrate and use $\|X-Y\|\leq\int\|\dot\gamma\|$; take the infimum
over paths. Apply Eq.~\eqref{eq:loss-product-lipschitz} for the loss bound.
\end{proof}

This is a product/quotient comparison on a spectral band. Comparing
$\|B-B'\|+\|A-A'\|$ without aligning or quotienting the gauge has no
such meaning: equivalent predictors can have arbitrarily distant factors.

\subsection{Uniform stability of the actual update map}
Consider two datasets differing only at index $j$, trained with the same
initial quotient state and the same random sample indices $I_t$. Write
$s_t$ for the full algorithmic state, including momentum or adaptive
statistics when present. Let $d_*(s,s')$ be a metric that dominates
$d_g(X,X')$. A first-order gradient method can use $s=X$ and $d_*=d_g$.

\begin{assumption}[State-map regularity for stability]
\label{ass:state-stability}
For a common example $z$, the actual state update $U_{t,z}$ obeys
\begin{equation}
 d_*(U_{t,z}s,U_{t,z}s')\leq(1+L_U\eta_t)d_*(s,s').
 \label{eq:update-lipschitz}
\end{equation}
Changing the example at the same state obeys
\begin{equation}
 d_*(U_{t,z}s,U_{t,z'}s)\leq2C_U\eta_t.
 \label{eq:sample-injection}
\end{equation}
The output loss is $C_\ell$-Lipschitz in this metric, uniformly in $z$.
For a pure quotient parameter, Lemma~\ref{lem:product-sensitive} permits
$C_\ell=L_\ell\sqrt{2b}$. These inequalities are assumptions on the
implemented update. They do not follow from loss smoothness alone for an
arbitrary optimizer with state.
\end{assumption}

\begin{theorem}[Nonconvex quotient-state stability]
\label{thm:stability}
Under Assumption~\ref{ass:state-stability}, one-example stochastic
sampling with $I_t$ uniform on $\{1,\ldots,n\}$ gives
\begin{align}
 \mathbb E d_*(s_T,s_T')
 &\leq\frac{2C_U}{n}\sum_{t<T}\eta_t
                  \prod_{k=t+1}^{T-1}(1+L_U\eta_k),
 \label{eq:stability-distance}\\
 \epsilon_{\rm stab}
 &\leq\frac{2C_\ell C_U}{n}\sum_{t<T}\eta_t
                  \prod_{k=t+1}^{T-1}(1+L_U\eta_k).
 \label{eq:stability-loss}
\end{align}
Here the supremum is over adjacent datasets and test examples, and the
expectation is over the common training randomness. In particular,
\begin{equation}
 \epsilon_{\rm stab}\leq\frac{2C_\ell C_U}{n}
    \exp\left(L_U\sum_{t<T}\eta_t\right)\sum_{t<T}\eta_t.
 \label{eq:nonconvex-stability-exp}
\end{equation}
The expected generalization gap satisfies
$|\mathbb E[R(X_T)-R_S(X_T)]|\leq\epsilon_{\rm stab}$.
\end{theorem}
\begin{proof}
If $I_t\ne j$, both updates use the same example, so
Eq.~\eqref{eq:update-lipschitz} applies directly. If $I_t=j$, insert an
intermediate update using the first dataset's example at $s_t'$.
The triangle inequality gives
\begin{equation}
 d_*(s_{t+1},s_{t+1}')\leq
 (1+L_U\eta_t)d_*(s_t,s_t')+2C_U\eta_t.
\end{equation}
The event has probability $1/n$. Conditional expectation therefore yields
\begin{equation}
 \mathbb E_t d_*(s_{t+1},s_{t+1}')\leq
 (1+L_U\eta_t)d_*(s_t,s_t')+2C_U\eta_t/n.
\end{equation}
Unrolling from zero initial distance proves
Eq.~\eqref{eq:stability-distance}. Loss Lipschitzness proves
Eq.~\eqref{eq:stability-loss}, and $1+x\leq e^x$ proves the simpler
exponential bound.

For completeness, take an independent replacement $z_i'$ for each $z_i$
and let $S^{(i)}$ be the dataset with only that replacement. Exchangeability
gives
\begin{align}
 \mathbb E[R(X_T(S))-R_S(X_T(S))]
 =\frac1n\sum_i\mathbb E\big[
 \ell(X_T(S^{(i)});z_i)-\ell(X_T(S);z_i)\big].
\end{align}
Each summand is bounded in absolute value by the uniform stability
estimate. This is the replace-one argument for stability-based
generalization \cite{bousquet2002stability,hardt2016train}.
\end{proof}

The exponential factor in Eq.~\eqref{eq:nonconvex-stability-exp} is
necessary in this general nonconvex argument. A claim of only
$O(\sum_t\eta_t/n)$ needs a nonexpansive update, convexity-like structure,
or another argument that controls expansion. Quotienting does not remove
all expansive directions of the task loss.

\begin{corollary}[Contractive stability ceiling]
\label{cor:contractive-stability}
If common-example updates instead contract with factor $1-\mu_U\eta$
at a constant step and the same sample-injection bound holds, then
\begin{equation}
 \epsilon_{\rm stab}\leq
 \frac{2C_\ell C_U}{\mu_Un}
       \{1-(1-\mu_U\eta)^T\}
 \leq\frac{2C_\ell C_U}{\mu_Un}.
 \label{eq:contractive-stability}
\end{equation}
\end{corollary}
\begin{proof}
Replace the multiplier in the preceding recursion by
$1-\mu_U\eta$ and sum the resulting geometric series.
\end{proof}

The contraction in this corollary concerns the stochastic per-example state
map. A PL inequality for the average training objective is not enough to
establish it. Proposition~\ref{prop:convex-counterexample} already shows
why even ambient product convexity can fail to imply the required
contraction for a factor-induced update.

\subsection{Why momentum requires a state metric}
To see the role of optimizer state explicitly, first consider a flat chart
with an exact transported identification and a constant step. For
$v^+=\beta v+(1-\beta)\nabla\ell_z(x)$, $x^+=x-\eta v^+$,
assume gradients have norm at most $G$ and are $L$-Lipschitz. Then
\begin{align}
 \delta_v^+&\leq\beta\delta_v+(1-\beta)L\delta_x,\\
 \delta_x^+&\leq\delta_x+\eta\delta_v^+.
\end{align}
For $0<\beta<1$, define the weighted state distance
\begin{equation}
 D=\delta_x+\frac{\eta\beta}{1-\beta}\delta_v.
 \label{eq:momentum-state-metric}
\end{equation}
Multiplication and collection of coefficients give
\begin{equation}
 D^+\leq(1+\eta L)D.
 \label{eq:momentum-state-lipschitz}
\end{equation}
Changing a single sampled gradient contributes at most $2\eta G$ to $D^+$:
the velocity injection is $2(1-\beta)G$ and its combined coefficient is
$\eta+\eta\beta/(1-\beta)=\eta/(1-\beta)$.
Thus Theorem~\ref{thm:stability} applies in this flat model with
$C_U=G$, rather than ignoring the velocity state.

On the quotient manifold one must add the sensitivity of the retraction,
the dependence of tangent identification on the base point, and the
transported-state norm. Smoothness on a compact regular chart makes the
associated constants finite, but they can depend on $a^{-1}$, the velocity
bound, and $\beta$. The appropriate statement is
Assumption~\ref{ass:state-stability} with these constants included.
For a time-varying stepsize the weight in
Eq.~\eqref{eq:momentum-state-metric} also changes; its change must be
controlled instead of reusing the constant-step calculation verbatim.
For Adam, the state additionally contains the elementwise second moment
and the regularizing denominator. A quotient stability theorem for that
state requires its own coordinate-consistent update definition.

\subsection{Function-space capacity from a nuclear-norm budget}
We next give a result that does not depend on the algorithm. Let a scalar
linear prediction use a matrix feature $Z(z)$ and product $X$:
$h_X(z)=\langle X,Z(z)\rangle$. Suppose
$\|Z(z)\|\leq R_Z$ and consider $\mathcal X_\tau=\{X:\|X\|_*\leq\tau\}$.
The frozen base matrix contributes only a fixed offset and does not affect
the centered Rademacher complexity.

\begin{theorem}[Gauge-invariant nuclear-norm generalization]
\label{thm:nuclear-generalization}
For the scalar predictor class just defined, its empirical Rademacher
complexity obeys
\begin{equation}
 \widehat{\mathfrak R}_n(\mathcal H_\tau)
 =\frac{\tau}{n}\mathbb E_\epsilon
       \left\|\sum_{i=1}^n\epsilon_i Z(z_i)\right\|_2
 \leq\frac{\tau R_Z}{\sqrt n}.
 \label{eq:nuclear-rademacher}
\end{equation}
If $\ell(h;z)\in[0,1]$ is $L_h$-Lipschitz in its scalar prediction,
then with probability at least $1-\delta$, simultaneously for this class,
\begin{equation}
 R(X)\leq R_S(X)+\frac{2L_h\tau R_Z}{\sqrt n}
                   +3\sqrt{\frac{\log(2/\delta)}{2n}}.
 \label{eq:nuclear-gen}
\end{equation}
For balanced factors, the budget is exactly
$\tau\geq\tfrac12(\|B\|^2+\|A\|^2)=\|X\|_*$.
For rank at most $r$ and $\|X\|_2\leq b$, one may take $\tau=rb$.
\end{theorem}
\begin{proof}
The dual norm of the nuclear norm is the spectral norm. Taking the
supremum over $\|X\|_*\leq\tau$ inside the Rademacher definition gives
the equality in Eq.~\eqref{eq:nuclear-rademacher}. The spectral norm is
at most the Frobenius norm, and Jensen's inequality and independence of
the signs give
\begin{equation}
 \mathbb E_\epsilon\left\|\sum_i\epsilon_i Z_i\right\|_2
 \leq\left(\mathbb E_\epsilon\left\|\sum_i\epsilon_i Z_i\right\|^2\right)^{1/2}
 =\left(\sum_i\|Z_i\|^2\right)^{1/2}\leq R_Z\sqrt n.
\end{equation}
The scalar contraction inequality bounds the loss complexity by $L_h$
times predictor complexity; centering the loss by its value at the fixed
zero predictor leaves that bound unchanged. The standard empirical
Rademacher uniform-deviation inequality for bounded losses gives
Eq.~\eqref{eq:nuclear-gen} \cite{bartlett2002rademacher}.
The factor identity is Theorem~\ref{thm:balanced-optimal}; the rank/spectral
bound follows by summing at most $r$ singular values each at most $b$.
\end{proof}

The capacity depends on the product budget, not on the energy of an
arbitrarily rescaled representative. Conversely, balancing a product with
a large nuclear norm does not shrink that nuclear norm. Therefore this
theorem does not establish a smaller function class for balance-only than
for an unbalanced method reaching exactly the same products.

\subsection{Intrinsic-dimension bounds in regular charts}
To expose the quotient dimension more directly, let
$\Psi_j:\Theta_j\subset\mathbb R^{D}\to\mathcal M_r$,
$D=r(m+d-r)$, be a finite atlas covering a compact product class
$\mathcal C$. Suppose each $\Theta_j$ lies in a Euclidean ball of radius
$R_j$ and
\begin{equation}
 \|\Psi_j(\theta)-\Psi_j(\theta')\|\leq C_j\|\theta-\theta'\|.
 \label{eq:chart-lipschitz}
\end{equation}
Such bounds exist after refining charts on a compact regular class;
their sizes and number can deteriorate when the class approaches a rank
boundary. Let $J$ be the number of charts, $C=\max C_j$, and
$R_\Theta=\max R_j$.

\begin{theorem}[Local quotient dimension in a uniform bound]
\label{thm:chart-generalization}
Under Eq.~\eqref{eq:loss-product-lipschitz}, for every $\varepsilon>0$,
with probability at least $1-\delta$,
\begin{equation}
 \sup_{X\in\mathcal C}|R(X)-R_S(X)|
 \leq2L_\ell\varepsilon+
 \sqrt{\frac{\log(2J/\delta)+
 D\log(1+2CR_\Theta/\varepsilon)}{2n}}.
 \label{eq:chart-generalization}
\end{equation}
\end{theorem}
\begin{proof}
A bounded subset of a $D$-dimensional radius-$R_\Theta$ Euclidean ball
has an $\varepsilon/C$ net with at most
$(1+2CR_\Theta/\varepsilon)^D$ elements, by the disjoint-ball volume
argument. Choose the centers within the chart domain. Mapping these nets
through Eq.~\eqref{eq:chart-lipschitz} gives an $\varepsilon$ net for
$\mathcal C$ with at most $J(1+2CR_\Theta/\varepsilon)^D$ elements.
Hoeffding's inequality and a union bound control the absolute
population--empirical difference on this finite net by the square-root
term. Approximating an arbitrary predictor by a net point changes both
its population and empirical losses by at most $L_\ell\varepsilon$,
which proves Eq.~\eqref{eq:chart-generalization}.
\end{proof}

This bound removes the $r^2$ redundant coordinates in the dimension of a
regular local parametrization, with explicit chart constants rather than
a claim that geometry provides $r^2$ fewer statistical samples for free.
A coarse global factor-based covering bound can use a larger exponent
and different constants; comparing those two upper bounds alone does not
prove a statistical advantage of one optimizer.

\subsection{Multilayer networks and spectral propagation}
Consider a depth-$H$ feed-forward predictor with weights
\begin{equation}
 W_l=W_{0,l}+c_lX_l,\qquad\operatorname{rank}X_l=r_l,
 \quad\|X_l\|_2\leq b_l,
\end{equation}
and fixed biases, or biases included in the parameter covers below.
Assume the activations are one-Lipschitz and map zero to zero, and
$\|x\|_2\leq R_x$. Define
\begin{equation}
 s_l=\|W_{0,l}\|_2+|c_l|b_l,
 \qquad A_l=R_x|c_l|\prod_{j\ne l}s_j.
 \label{eq:network-spectral-constants}
\end{equation}
For activations with Lipschitz constants $\rho_l$, insert the corresponding
products of $\rho_l$ in these constants. In particular a frequency-scaled
sine activation cannot be silently treated as one-Lipschitz.

\begin{lemma}[Telescoping sensitivity of a product-parametrized network]
\label{lem:network-sensitivity}
For two networks in the stated class with the same frozen weights,
\begin{equation}
 \|h_X(x)-h_Y(x)\|_2
       \leq\sum_{l=1}^H A_l\|X_l-Y_l\|_2.
 \label{eq:network-telescope}
\end{equation}
If a layer also multiplies a fixed propagation operator $P_l$, its spectral
norm contributes a factor $\|P_l\|_2$ in the appropriate propagation
products. Symmetric degree-normalized adjacency with self-loops has
spectral norm at most one, so it does not enlarge those products.
\end{lemma}
\begin{proof}
Replace the layers one at a time and telescope the output differences.
Before layer $l$, the input norm is at most
$R_x\prod_{j<l}s_j$. The weight replacement changes that layer's
preactivation by at most $|c_l|\|X_l-Y_l\|_2$ times the input norm.
The remaining layers amplify it by at most $\prod_{j>l}s_j$ because
each activation is one-Lipschitz. Sum over $l$. Fixed propagation
operators are bounded by submultiplicativity. The normalized adjacency
is similar to a stochastic matrix and is symmetric, so all its
eigenvalues lie in $[-1,1]$.
\end{proof}

\begin{theorem}[A network capacity bound with quotient coordinates]
\label{thm:network-generalization}
Suppose layer $l$ has a cover satisfying Theorem~\ref{thm:chart-generalization}
with dimension $D_l=r_l(m_l+d_l-r_l)$, chart count $J_l$, and constants
$C_l,R_l$. The scalar loss is $L_{\rm out}$-Lipschitz in the network
output and lies in $[0,1]$. Then, with probability at least $1-\delta$,
\begin{align}
 \sup_X|R(X)-R_S(X)|
 &\leq2L_{\rm out}\sum_l A_l\varepsilon_l\notag\\
 &\quad+\sqrt{\frac{\log(2/\delta)+\sum_l\log J_l+
 \sum_lD_l\log(1+2C_lR_l/\varepsilon_l)}{2n}}.
 \label{eq:network-generalization}
\end{align}
Additional trainable dense layers, embeddings, attention parameters, and
biases require their own covering terms; they cannot be omitted because
only the adapters are called the method's parameters.
\end{theorem}
\begin{proof}
Take the Cartesian product of layerwise nets. Their cardinalities
multiply, so their logarithms add. By Lemma~\ref{lem:network-sensitivity},
each resulting output approximation has norm error at most
$\sum_l A_l\varepsilon_l$, using $\|\cdot\|_2\leq\|\cdot\|$ for the
layer net errors. Apply Hoeffding and a union bound on this finite product
net, then pay the output-loss Lipschitz approximation cost twice.
\end{proof}

These network formulas concern the stated feed-forward class. Transformer
attention contains a data-dependent softmax operator; its Lipschitz and
norm constants must be bounded as part of a corresponding full-network
cover. The theorem is not a shortcut to a dimension-free Transformer
generalization guarantee. Similarly, transductive nodes in a single graph
are not independent graph observations. Applying the IID concentration
steps to Cora, Citeseer, PubMed, or ogbn-arxiv requires an appropriate
sampling model or a separate transductive bound. Sequential language tokens
also do not constitute IID examples. We use held-out task metrics for
those experiments and do not present them as empirical confirmations of
the IID theorems.

\subsection{Model selection across the full proposal grid}
Let $M$ predictors be trained without using a validation sample of size
$n_v$. They may include every successful hyperparameter configuration,
checkpoint, rank, and architecture under consideration. Conditional on the
training data and algorithmic randomness, the predictors are fixed.

\begin{proposition}[Validation selection cost]
\label{prop:selection}
For bounded validation losses, with probability at least $1-\delta$,
the predictor $\widehat h$ minimizing validation loss among the $M$
candidates obeys
\begin{equation}
 R(\widehat h)\leq\min_{1\leq j\leq M}R(h_j)
           +2\sqrt{\frac{\log(2M/\delta)}{2n_v}}.
 \label{eq:validation-selection}
\end{equation}
\end{proposition}
\begin{proof}
Hoeffding and a union bound yield simultaneous deviations at most
$a=\sqrt{\log(2M/\delta)/(2n_v)}$. Thus
$R(\widehat h)\leq\widehat R_v(\widehat h)+a
\leq\widehat R_v(h_\star)+a\leq R(h_\star)+2a$,
where $h_\star$ is the best population-risk candidate.
\end{proof}

The candidate count includes searched checkpoints when validation selects
an epoch. Replicating configurations whose damping or balance interval
does not affect a baseline is an execution alias, not independent evidence.
If the grid itself is adaptively redesigned using validation feedback,
the fixed-candidate conditional argument no longer applies without
accounting for that adaptivity. Test data are not used to pick the grid's
winning method or checkpoint. Reporting all attempted tasks and failed
configurations is necessary for interpreting a broad search for positive
results.

\subsection{Combining generalization with actual optimization guarantees}
\begin{theorem}[Population excess risk under explicit empirical optimization]
\label{thm:excess-risk}
Let $\mathcal C$ be a fixed predictor class containing the output $X_T$.
If a uniform deviation bound
$\sup_{X\in\mathcal C}|R(X)-R_S(X)|\leq U_n$ holds and the algorithm
achieves
\begin{equation}
 R_S(X_T)-\inf_{X\in\mathcal C}R_S(X)\leq\varepsilon_{\rm opt},
\end{equation}
then on that event
\begin{equation}
 R(X_T)-\inf_{X\in\mathcal C}R(X)
           \leq\varepsilon_{\rm opt}+2U_n.
 \label{eq:excess-risk-uniform}
\end{equation}
Alternatively, if the algorithm has expected uniform stability
$\epsilon_{\rm stab}$ and expected empirical suboptimality at most
$\varepsilon_{\rm opt}$, then
\begin{equation}
 \mathbb E[R(X_T)]-\inf_{X\in\mathcal C}R(X)
           \leq\varepsilon_{\rm opt}+\epsilon_{\rm stab}.
 \label{eq:excess-risk-stable}
\end{equation}
The statements allow infima by taking limits; attainment is unnecessary.
\end{theorem}
\begin{proof}
For any $X\in\mathcal C$, add and subtract empirical risks:
\begin{align}
 R(X_T)-R(X)
 &=[R(X_T)-R_S(X_T)]+[R_S(X_T)-R_S(X)]
                         +[R_S(X)-R(X)]\\
 &\leq2U_n+\varepsilon_{\rm opt}.
\end{align}
Take an infimizing sequence to prove Eq.~\eqref{eq:excess-risk-uniform}.
For the expectation bound, compare to any fixed $X$ independent of $S$.
Expected stability bounds the first bracket by $\epsilon_{\rm stab}$,
and the expected third bracket is zero by IID sampling. The empirical
optimization guarantee bounds the middle bracket. Take the infimum over
fixed comparison points.
\end{proof}

If the empirical objective has the quotient PL property over an invariant
optimization region containing its class minimum, Theorem~\ref{thm:pl}
provides $\varepsilon_{\rm opt}\leq(1-\mu_g\eta)^T\Delta_0$ in the
deterministic case. Corollary~\ref{cor:stochastic-pl} supplies the
corresponding noisy optimization floor. In contrast, the nonconvex
stationarity bounds in Theorems~\ref{thm:stochastic} and
\ref{thm:momentum} do not give empirical suboptimality without an
additional landscape assumption. Substituting a squared-gradient bound
for $\varepsilon_{\rm opt}$ would be invalid.

Any approximation error from restricting to low rank is separate:
\begin{equation}
 R(X_T)-\inf_{h\in\mathcal H}R(h)
 =\big[R(X_T)-\inf_{X\in\mathcal C}R(X)\big]
  +\big[\inf_{X\in\mathcal C}R(X)-\inf_{h\in\mathcal H}R(h)\big].
\end{equation}
Increasing rank can reduce this approximation term while enlarging capacity
and computational cost. Quotient optimization changes how the rank-$r$
class is traversed; it does not remove that tradeoff.

\subsection{Consequences for the empirical claims}
The theory supports monitoring a coherent list of quantities: product rank
margin, actual tangent-gradient norm, update-map sensitivity, represented
spectral and nuclear norms, training loss, and validation-selected task
performance. It does not support declaring improved generalization from
factor balancing alone. At a fixed product, balancing preserves the entire
predictor, including its training and population risks. If a SIREN layer
computes its factors only through $BA$ before applying a sine, balancing
also preserves that layer exactly; a claim that it removes a useful
factor-scale-dependent representation is false for that architecture.
Different results after training can reflect different product-space
optimization paths, frequency-sensitive conditioning, optimizer state, or
implementation error, and must be investigated as such.

The generalization analysis is therefore both constructive and scoped:
it gives explicit bounds when their statistical and optimization assumptions
hold, identifies the intrinsic rank-$r$ dimension in regular coordinates,
and explains why gauge-invariant quantities are the appropriate objects.
It leaves real task superiority and the sign of a train--test gap difference
to completed, properly controlled experiments.
